\Chapter{6}

\Verse{1}
{
    \zA{却说秦氏因听见宝玉梦中唤他的乳名，心中纳闷，又不好细问。}
}
{
    \eA{Mrs. Ch'in, to resume our narrative, upon hearing Pao-yü call her in his
        dream by her infant name, was at heart very exercised, but she did not
        however feel at liberty to make any minute inquiry.}
}
{
}

\Verse{2}
{
    \zA{彼时宝玉迷迷 惑惑，若有所失，遂起身解怀整衣。}
}
{
    \eA{Pao-yü was, at this time, in such a dazed state, as if he had lost
        something, and the servants promptly gave him a decoction of lungngan.}
}
{
}

\Verse{3}
{
    \zA{袭人过来给他系裤带时，刚伸手至大腿处，只 觉冰冷粘湿的一片，吓的忙褪回手来，问：“是怎么了？”}
}
{
    \eA{After he had taken a few sips, he forthwith rose and tidied his clothes. Hsi
        Jen put out her hand to fasten the band of his garment, and as soon as she
        did so, and it came in contact with his person,}
}
{
}

\Verse{4}
{
    \zA{宝玉红了脸，把他的手 一捻。}
}
{
    \eA{it felt so icy cold to the touch, covered as it was all over with
        perspiration,}
}
{
}

\Verse{5}
{
    \zA{袭人本是个聪明女子，年纪又比宝玉大两岁，近来也渐省人事。}
}
{
    \eA{that she speedily withdrew her hand in utter surprise. "What's the matter
        with you?" she exclaimed. A blush suffused Pao-yü's face, and he took Hsi
        Jen's hand in a tight grip.}
}
{
}

\Verse{6}
{
    \zA{今见宝玉如 此光景，心中便觉察了一半，不觉把个粉脸羞的飞红，遂不好再问。}
}
{
    \eA{Hsi Jen was a girl with all her wits about her; she was besides a couple of
        years older than Pao-yü and had recently come to know something of the
        world, so that at the sight of his state, she to a great extent readily
        accounted for the reason in her heart. From modest shame, she unconsciously
        became purple in the face,}
}
{
}

\Verse{7}
{
    \zA{仍旧理好衣裳， 随至贾母处来，胡乱吃过晚饭，过这边来，趁众奶娘丫鬟不在旁时，另取出一件中 衣与宝玉换上。}
}
{
    \eA{and not venturing to ask another question she continued adjusting his
        clothes. This task accomplished, she followed him over to old lady Chia's
        apartments; and after a hurry-scurry meal, they came back to this side, and
        Hsi Jen availed herself of the absence of the nurses and waiting-maids to
        hand Pao-yü another garment to change.}
}
{
}

\Verse{8}
{
    \zA{宝玉含羞央告道：“好姐姐，千万别告诉人。”}
}
{
    \eA{"Please, dear Hsi Jen, don't tell any one," entreated Pao-yü, with concealed
        shame. "What did you dream of?"}
}
{
}

\Verse{9}
{
    \zA{袭人也含着羞悄悄 的笑问道：“你为什么――”说到这里，把眼又往四下里瞧了瞧，才又问道：“那 是那里流出来的？”}
}
{
    \eA{inquired Hsi Jen, smiling, as she tried to stifle her blushes, "and whence
        comes all this perspiration?" "It's a long story," said Pao-yü, "which only
        a few words will not suffice to explain." He accordingly recounted minutely,
        for her benefit, the subject of his dream.}
}
{
}

\Verse{10}
{
    \zA{宝玉只管红着脸不言语，袭人却只瞅着他笑。}
}
{
    \eA{When he came to where the Fairy had explained to him the mysteries of love,
        Hsi Jen was overpowered with modesty and covered her face with her hands;}
}
{
}

\Verse{11}
{
    \zA{迟了一会，宝玉 才把梦中之事细说与袭人听。}
}
{
    \eA{and as she bent down, she gave way to a fit of laughter. Pao-yü had always
        been fond of Hsi Jen, on account of her gentleness,}
}
{
}

\Verse{12}
{
    \zA{说到云雨私情，羞的袭人掩面伏身而笑。}
}
{
    \eA{pretty looks and graceful and elegant manner, and he forthwith expounded to
        her all the mysteries he had been taught by the Fairy. Hsi Jen was, of
        course, well aware that dowager lady Chia had given her over to Pao-yü,}
}
{
}

\Verse{13}
{
    \zA{宝玉亦素喜 袭人柔媚姣俏，遂强拉袭人同领警幻所训之事，袭人自知贾母曾将他给了宝玉，也 无可推托的，扭捏了半日，无奈何，只得和宝玉温存了一番。}
}
{
    \eA{so that her present behaviour was likewise no transgression. And
        subsequently she secretly attempted with Pao-yü a violent flirtation, and
        lucky enough no one broke in upon them during their tête-à-tête. From this
        date, Pao-yü treated Hsi Jen with special regard, far more than he showed to
        the other girls,}
}
{
}

\Verse{14}
{
    \zA{自此宝玉视袭人更自 不同，袭人待宝玉也越发尽职了。}
}
{
    \eA{while Hsi Jen herself was still more demonstrative in her attentions to
        Pao-yü. But for a time we will make no further remark about them. As regards
        the household of the Jung mansion, the inmates may, on adding up the total
        number, not have been found many;}
}
{
}

\Verse{15}
{
    \zA{这话暂且不提。}
}
{
    \eA{yet, counting the high as well as the low,}
}
{
}

\Verse{16}
{
    \zA{且说荣府中合算起来，从上至下，也有三百馀口人，一天也有一二十件事，竟 如乱麻一般，没个头绪可作纲领。}
}
{
    \eA{there were three hundred persons and more. Their affairs may not have been
        very numerous, still there were, every day, ten and twenty matters to
        settle; in fact, the household resembled, in every way, ravelled hemp,
        devoid even of a clue-end,}
}
{
}

\Verse{17}
{
    \zA{正思从那一件事那一个人写起方妙，却好忽从千 里之外，芥豆之微，小小一个人家，因与荣府略有些瓜葛，这日正往荣府中来，因 此便就这一家说起，倒还是个头绪。}
}
{
    \eA{which could be used as an introduction. Just as we were considering what
        matter and what person it would be best to begin writing of, by a lucky
        coincidence suddenly from a distance of a thousand li, a person small and
        insignificant as a grain of mustard seed happened, on account of her distant
        relationship with the Jung family, to come on this very day to the Jung
        mansion on a visit.}
}
{
}

\Verse{18}
{
    \zA{原来这小小之家，姓王，乃本地人氏，祖上也做过一个小小京官，昔年曾与凤 姐之祖王夫人之父认识。}
}
{
    \eA{We shall therefore readily commence by speaking of this family, as it after
        all affords an excellent clue for a beginning. The surname of this mean and
        humble family was in point of fact Wang. They were natives of this district.
        Their ancestor had filled a minor office in the capital,}
}
{
}

\Verse{19}
{
    \zA{因贪王家的势利，便连了宗，认作侄儿。}
}
{
    \eA{and had, in years gone by, been acquainted with lady Feng's grandfather,
        that is madame Wang's father.}
}
{
}

\Verse{20}
{
    \zA{那时只有王夫人 之大兄凤姐之父与王夫人随在京的知有此一门远族，馀者也皆不知。}
}
{
    \eA{Being covetous of the influence and affluence of the Wang family, he
        consequently joined ancestors with them, and was recognised by them as a
        nephew. At that time, there were only madame Wang's eldest brother, that is
        lady Feng's father,}
}
{
}

\Verse{21}
{
    \zA{目今其祖早 故，只有一个儿子，名唤王成，因家业萧条，仍搬出城外乡村中住了。}
}
{
    \eA{and madame Wang herself, who knew anything of these distant relations, from
        the fact of having followed their parents to the capital. The rest of the
        family had one and all no idea about them. This ancestor had, at this date,}
}
{
}

\Verse{22}
{
    \zA{王成亦相继 身故，有子小名狗儿，娶妻刘氏，生子小名板儿；又生一女，名唤青儿：一家四口， 以务农为业。}
}
{
    \eA{been dead long ago, leaving only one son called Wang Ch'eng. As the family
        estate was in a state of ruin, he once more moved outside the city walls and
        settled down in his native village. Wang Ch'eng also died soon after his
        father, leaving a son, known in his infancy as Kou Erh,}
}
{
}

\Verse{23}
{
    \zA{因狗儿白日间自作些生计，刘氏又操井臼等事，青板姊弟两个无人照 管，狗儿遂将岳母刘老老接来，一处过活。}
}
{
    \eA{who married a Miss Liu, by whom he had a son called by the infant name of
        Pan Erh, as well as a daughter, Ch'ing Erh. His family consisted of four,
        and he earned a living from farming. As Kou Erh was always busy with
        something or other during the day and his wife, dame Liu, on the other hand,}
}
{
}

\Verse{24}
{
    \zA{这刘老老乃是个久经世代的老寡妇，膝 下又无子息，只靠两亩薄田度日。}
}
{
    \eA{drew the water, pounded the rice and attended to all the other domestic
        concerns, the brother and sister, Ch'ing Erh and Pan Erh, the two of them,
        had no one to look after them. (Hence it was that) Kou Erh brought over his
        mother-in-law,}
}
{
}

\Verse{25}
{
    \zA{如今女婿接了养活。}
}
{
    \eA{old goody Liu, to live with them. This goody Liu was an old widow,}
}
{
}

\Verse{26}
{
    \zA{岂不愿意呢，遂一心一计， 帮着女儿女婿过活。}
}
{
    \eA{with a good deal of experience. She had besides no son round her knees, so
        that she was dependent for her maintenance on a couple of acres of poor
        land, with the result that when her son-in-law received her in his home,}
}
{
}

\Verse{27}
{
    \zA{因这年秋尽冬初，天气冷将上来，家中冬事未办，狗儿未免心中烦躁，吃了几 杯闷酒，在家里闲寻气恼，刘氏不敢顶撞。}
}
{
    \eA{she naturally was ever willing to exert heart and mind to help her daughter
        and her son-in-law to earn their living. This year, the autumn had come to
        an end, winter had commenced, and the weather had begun to be quite cold. No
        provision had been made in the household for the winter months, and Kou Erh
        was, inevitably,}
}
{
}

\Verse{28}
{
    \zA{因此刘老老看不过，便劝道：“姑爷， 你别嗔着我多嘴：咱们村庄人家儿，那一个不是老老实实，守着多大碗儿吃多大的
        饭呢!你皆因年小时候，托着老子娘的福，吃喝惯了，如今所以有了钱就顾头不顾 尾，没了钱就瞎生气，成了什么男子汉大丈夫了!如今咱们虽离城住着，终是天子
        脚下。}
}
{
    \eA{exceedingly exercised in his heart. Having had several cups of wine to
        dispel his distress, he sat at home and tried to seize upon every trifle to
        give vent to his displeasure. His wife had not the courage to force herself
        in his way, and hence goody Liu it was who encouraged him, as she could not
        bear to see the state of the domestic affairs. "Don't pull me up for talking
        too much," she said; "but who of us country people isn't honest and
        open-hearted?}
}
{
}

\Verse{29}
{
    \zA{这长安城中遍地皆是钱，只可惜没人会去拿罢了。}
}
{
    \eA{As the size of the bowl we hold, so is the quantity of the rice we eat. In
        your young days, you were dependent on the support of your old father,}
}
{
}

\Verse{30}
{
    \zA{在家跳蹋也没用！”}
}
{
    \eA{so that eating and drinking became quite a habit with you;}
}
{
}

\Verse{31}
{
    \zA{狗儿 听了道：“你老只会在炕头上坐着混说，难道叫我打劫去不成？”}
}
{
    \eA{that's how, at the present time, your resources are quite uncertain; when
        you had money, you looked ahead, and didn't mind behind; and now that you
        have no money, you blindly fly into huffs.}
}
{
}

\Verse{32}
{
    \zA{刘老老说道：“谁 叫你去打劫呢？}
}
{
    \eA{A fine fellow and a capital hero you have made! Living though we now be away
        from the capital,}
}
{
}

\Verse{33}
{
    \zA{也到底大家想个方法儿才好。}
}
{
    \eA{we are after all at the feet of the Emperor; this city of Ch'ang Ngan is
        strewn all over with money,}
}
{
}

\Verse{34}
{
    \zA{不然那银子钱会自己跑到咱们家里来 不成？”}
}
{
    \eA{but the pity is that there's no one able to go and fetch it away; and it's
        no use your staying at home and kicking your feet about." "All you old lady
        know,"}
}
{
}

\Verse{35}
{
    \zA{狗儿冷笑道：“有法儿还等到这会子呢!我又没有收税的亲戚、做官的朋 友，有什么法子可想的？}
}
{
    \eA{rejoined Kou Erh, after he had heard what she had to say, "is to sit on the
        couch and talk trash! Is it likely you would have me go and play the
        robber?" "Who tells you to become a robber?" asked goody Liu. "But it would
        be well, after all, that we should put our heads together and devise some
        means;}
}
{
}

\Verse{36}
{
    \zA{就有，也只怕他们未必来理我们呢。”}
}
{
    \eA{for otherwise, is the money, pray, able of itself to run into our house?"
        "Had there been a way," observed Kou Erh,}
}
{
}

\Verse{37}
{
    \zA{刘老老道“这倒也 不然。}
}
{
    \eA{smiling sarcastically, "would I have waited up to this moment?}
}
{
}

\Verse{38}
{
    \zA{‘谋事在人，成事在天’，咱们谋到了，靠菩萨的保佑，有些机会，也未可 知。}
}
{
    \eA{I have besides no revenue collectors as relatives, or friends in official
        positions; and what way could we devise? 'But even had I any, they wouldn't
        be likely, I fear, to pay any heed to such as ourselves!"}
}
{
}

\Verse{39}
{
    \zA{我倒替你们想出一个机会来。}
}
{
    \eA{"That, too, doesn't follow," remarked goody Liu; "the planning of affairs
        rests with man,}
}
{
}

\Verse{40}
{
    \zA{当日你们原是和金陵王家连过宗的。}
}
{
    \eA{but the accomplishment of them rests with Heaven. After we have laid our
        plans, we may, who can say, by relying on the sustenance of the gods,}
}
{
}

\Verse{41}
{
    \zA{二十年前， 他们看承你们还好，如今是你们拉硬屎，不肯去就和他，才疏远起来。}
}
{
    \eA{find some favourable occasion. Leave it to me, I'll try and devise some
        lucky chance for you people! In years gone by, you joined ancestors with the
        Wang family of Chin Ling, and twenty years back, they treated you with
        consideration;}
}
{
}

\Verse{42}
{
    \zA{想当初我和 女儿还去过一遭，他家的二小姐着实爽快会待人的，倒不拿大，如今现是荣国府贾 二老爷的夫人。}
}
{
    \eA{but of late, you've been so high and mighty, and not condescended to go and
        bow to them, that an estrangement has arisen. I remember how in years gone
        by, I and my daughter paid them a visit. The second daughter of the family
        was really so pleasant and knew so well how to treat people with kindness,}
}
{
}

\Verse{43}
{
    \zA{听见他们说，如今上了年纪，越发怜贫恤老的了，又爱斋僧布施。}
}
{
    \eA{and without in fact any high airs! She's at present the wife of Mr. Chia,
        the second son of the Jung Kuo mansion; and I hear people say that now that
        she's advanced in years,}
}
{
}

\Verse{44}
{
    \zA{如今王府虽升了官儿，只怕二姑太太还认的咱们，你为什么不走动走动？}
}
{
    \eA{she's still more considerate to the poor, regardful of the old, and very
        fond of preparing vegetable food for the bonzes and performing charitable
        deeds. The head of the Wang mansion has,}
}
{
}

\Verse{45}
{
    \zA{或者他还 念旧，有些好处也未可知。}
}
{
    \eA{it is true, been raised to some office on the frontier, but I hope that this
        lady Secunda will anyhow notice us.}
}
{
}

\Verse{46}
{
    \zA{只要他发点好心，拔根寒毛，比咱们的腰还壮呢。”}
}
{
    \eA{How is it then that you don't find your way as far as there; for she may
        possibly remember old times, and some good may, no one can say,}
}
{
}

\Verse{47}
{
    \zA{刘 氏接口道：“你老说的好，你我这样嘴脸，怎么好到他门上去？}
}
{
    \eA{come of it? I only wish that she would display some of her kind-heartedness,
        and pluck one hair from her person which would be, yea thicker than our
        waist." "What you suggest, mother,}
}
{
}

\Verse{48}
{
    \zA{只怕他那门上人也 不肯进去告诉，没的白打嘴现世的！”}
}
{
    \eA{is quite correct," interposed Mrs. Liu, Kou Erh's wife, who stood by and
        took up the conversation, "but with such mouth and phiz as yours and mine,}
}
{
}

\Verse{49}
{
    \zA{谁知狗儿利名心重，听如此说，心下便有些活动；又听他妻子这番话，便笑道： “老老既这么说，况且当日你又见过这姑太太一次，为什么不你老人家明日就去走
        一遭，先试试风头儿去？”}
}
{
    \eA{how could we present ourselves before her door? Why I fear that the man at
        her gate won't also like to go and announce us! and we'd better not go and
        have our mouths slapped in public!" Kou Erh, who would have thought it,
        prized highly both affluence and fame, so that when he heard these remarks,
        he forthwith began to feel at heart a little more at ease.}
}
{
}

\Verse{50}
{
    \zA{刘老老道：“哎哟!可是说的了：‘侯门似海。’}
}
{
    \eA{When he furthermore heard what his wife had to say, he at once caught up the
        word as he smiled. "Old mother," he rejoined;}
}
{
}

\Verse{51}
{
    \zA{我是 个什么东西儿!他家人又不认得我，去了也是白跑。”}
}
{
    \eA{"since that be your idea, and what's more, you have in days gone by seen
        this lady on one occasion, why shouldn't you, old lady, start to-morrow on a
        visit to her and first ascertain how the wind blows!"}
}
{
}

\Verse{52}
{
    \zA{狗儿道：“不妨，我教给你 个法儿。}
}
{
    \eA{"Ai Ya!" exclaimed old Goody, "It may very well be said that the marquis'
        door is like the wide ocean!}
}
{
}

\Verse{53}
{
    \zA{你竟带了小板儿先去找陪房周大爷，要见了他，就有些意思了。}
}
{
    \eA{what sort of thing am I? why the servants of that family wouldn't even
        recognise me! even were I to go, it would be on a wild goose chase." "No
        matter about that," observed Kou Erh;}
}
{
}

\Verse{54}
{
    \zA{这周大爷 先时和我父亲交过一桩事，我们本极好的。”}
}
{
    \eA{"I'll tell you a good way; you just take along with you, your grandson,
        little Pan Erh, and go first and call upon Chou Jui, who is attached to that
        household;}
}
{
}

\Verse{55}
{
    \zA{刘老老道：“我也知道。}
}
{
    \eA{and when once you've seen him, there will be some little chance.}
}
{
}

\Verse{56}
{
    \zA{只是许多时 不走动，知道他如今是怎样？}
}
{
    \eA{This Chou Jui, at one time, was connected with my father in some affair or
        other, and we were on excellent terms with him."}
}
{
}

\Verse{57}
{
    \zA{――这也说不得了!你又是个男人，这么个嘴脸，自然 去不得；我们姑娘年轻的媳妇儿，也难卖头卖脚的。}
}
{
    \eA{"That I too know," replied goody Liu, "but the thing is that you've had no
        dealings with him for so long, that who knows how he's disposed towards us
        now? this would be hard to say. Besides, you're a man, and with a mouth and
        phiz like that of yours, you couldn't, on any account,}
}
{
}

\Verse{58}
{
    \zA{倒还是舍着我这副老脸去碰碰， 果然有好处，大家也有益。”}
}
{
    \eA{go on this errand. My daughter is a young woman, and she too couldn't very
        well go and expose herself to public gaze. But by my sacrificing this old
        face of mine,}
}
{
}

\Verse{59}
{
    \zA{当晚计议已定。}
}
{
    \eA{and by going and knocking it (against the wall) there may,}
}
{
}

\Verse{60}
{
    \zA{次日天未明时，刘老老便起来梳洗了。}
}
{
    \eA{after all, be some benefit and all of us might reap profit." That very same
        evening, they laid their plans,}
}
{
}

\Verse{61}
{
    \zA{又将板儿教了几句话。}
}
{
    \eA{and the next morning before the break of day, old goody Liu speedily got up,}
}
{
}

\Verse{62}
{
    \zA{五六岁的孩子， 听见带了他进城逛去，喜欢的无不应承。}
}
{
    \eA{and having performed her toilette, she gave a few useful hints to Pan Erh;
        who, being a child of five or six years of age, was, when he heard that he
        was to be taken into the city,}
}
{
}

\Verse{63}
{
    \zA{于是刘老老带了板儿，进城至宁荣街来。}
}
{
    \eA{at once so delighted that there was nothing that he would not agree to.
        Without further delay, goody Liu led off Pan Erh,}
}
{
}

\Verse{64}
{
    \zA{到了荣府大门前石狮子旁边，只见满门口的轿马。}
}
{
    \eA{and entered the city, and reaching the Ning Jung street, she came to the
        main entrance of the Jung mansion, where, next to the marble lions,}
}
{
}

\Verse{65}
{
    \zA{刘老老不敢过去，掸掸衣服，又 教了板儿几句话，然后溜到角门前，只见几个挺胸叠肚、指手画脚的人坐在大门上， 说东谈西的。}
}
{
    \eA{were to be seen a crowd of chairs and horses. Goody Liu could not however
        muster the courage to go by, but having shaken her clothes, and said a few
        more seasonable words to Pan Erh, she subsequently squatted in front of the
        side gate, whence she could see a number of servants,}
}
{
}

\Verse{66}
{
    \zA{刘老老只得蹭上来问：“太爷们纳福。”}
}
{
    \eA{swelling out their chests, pushing out their stomachs, gesticulating with
        their hands and kicking their feet about,}
}
{
}

\Verse{67}
{
    \zA{众人打量了一会，便问：“是 那里来的？”}
}
{
    \eA{while they were seated at the main entrance chattering about one thing and
        another. Goody Liu felt constrained to edge herself forward.}
}
{
}

\Verse{68}
{
    \zA{刘老老陪笑道：“我找太太的陪房周大爷的。}
}
{
    \eA{"Gentlemen," she ventured, "may happiness betide you!" The whole company of
        servants scrutinised her for a time.}
}
{
}

\Verse{69}
{
    \zA{烦那位太爷替我请他出 来。”}
}
{
    \eA{"Where do you come from?" they at length inquired. "I've come to look up Mr.
        Chou,}
}
{
}

\Verse{70}
{
    \zA{那些人听了，都不理他，半日方说道：“你远远的那墙畸角儿等着，一会子 他们家里就有人出来。”}
}
{
    \eA{an attendant of my lady's," remarked goody Liu, as she forced a smile;
        "which of you, gentlemen, shall I trouble to do me the favour of asking him
        to come out?" The servants, after hearing what she had to say, paid, the
        whole number of them,}
}
{
}

\Verse{71}
{
    \zA{内中有个年老的说道：“何苦误他的事呢？”}
}
{
    \eA{no heed to her; and it was after the lapse of a considerable time that they
        suggested: "Go and wait at a distance,}
}
{
}

\Verse{72}
{
    \zA{因向刘老老 道：“周大爷往南边去了。}
}
{
    \eA{at the foot of that wall; and in a short while, the visitors, who are in
        their house, will be coming out."}
}
{
}

\Verse{73}
{
    \zA{他在后一带住着，他们奶奶儿倒在家呢。}
}
{
    \eA{Among the party of attendants was an old man, who interposed, "Don't baffle
        her object," he expostulated; "why make a fool of her?"}
}
{
}

\Verse{74}
{
    \zA{你打这边绕到 后街门上找就是了。”}
}
{
    \eA{and turning to goody Liu: "This Mr. Chou," he said, "is gone south: his
        house is at the back row;}
}
{
}

\Verse{75}
{
    \zA{刘老老谢了，遂领着板儿绕至后门上，只见门上歇着些生意 担子，也有卖吃的，也有卖玩耍的，闹吵吵三二十个孩子在那里。}
}
{
    \eA{his wife is anyhow at home; so go round this way, until you reach the door,
        at the back street, where, if you will ask about her, you will be on the
        right track." Goody Liu, having expressed her thanks,}
}
{
}

\Verse{76}
{
    \zA{刘老老便拉住一 个道：“我问哥儿一声：有个周大娘在家么？”}
}
{
    \eA{forthwith went, leading Pan Erh by the hand, round to the back door, where
        she saw several pedlars resting their burdens. There were also those who
        sold things to eat,}
}
{
}

\Verse{77}
{
    \zA{那孩子翻眼瞅着道：“那个周大娘？}
}
{
    \eA{and those who sold playthings and toys; and besides these, twenty or thirty
        boys bawled and shouted,}
}
{
}

\Verse{78}
{
    \zA{我们这里周大娘有几个呢，不知那一个行当儿上的？”}
}
{
    \eA{making quite a noise. Goody Liu readily caught hold of one of them. "I'd
        like to ask you just a word, my young friend," she observed; "there's a Mrs.
        Chou here;}
}
{
}

\Verse{79}
{
    \zA{刘老老道：“他是太太的陪 房。”}
}
{
    \eA{is she at home?" "Which Mrs. Chou?" inquired the boy; "we here have three
        Mrs. Chous;}
}
{
}

\Verse{80}
{
    \zA{那孩子道：“这个容易，你跟了我来。”}
}
{
    \eA{and there are also two young married ladies of the name of Chou. What are
        the duties of the one you want,}
}
{
}

\Verse{81}
{
    \zA{引着刘老老进了后院，到一个院子 墙边，指道：“这就是他家。”}
}
{
    \eA{I wonder ?" "She's a waiting-woman of my lady," replied goody Liu. "It's
        easy to get at her," added the boy; "just come along with me." Leading the
        way for goody Liu into the backyard,}
}
{
}

\Verse{82}
{
    \zA{又叫道：“周大妈，有个老奶奶子找你呢。”}
}
{
    \eA{they reached the wall of a court, when he pointed and said, "This is her
        house.--Mother Chou!"}
}
{
}

\Verse{83}
{
    \zA{周瑞家的在内忙迎出来，问：“是那位？”}
}
{
    \eA{he went on to shout with alacrity; "there's an old lady who wants to see
        you." Chou Jui's wife was at home,}
}
{
}

\Verse{84}
{
    \zA{刘老老迎上来笑问道：“好啊？}
}
{
    \eA{and with all haste she came out to greet her visitor. "Who is it?" she
        asked. Goody Liu advanced up to her.}
}
{
}

\Verse{85}
{
    \zA{周 嫂子。”}
}
{
    \eA{"How are you," she inquired,}
}
{
}

\Verse{86}
{
    \zA{周瑞家的认了半日，方笑道：“刘老老，你好？}
}
{
    \eA{"Mrs. Chou?" Mrs. Chou looked at her for some time before she at length
        smiled and replied, "Old goody Liu, are you well? How many years is it since
        we've seen each other;}
}
{
}

\Verse{87}
{
    \zA{你说么，这几年不见，我 就忘了。}
}
{
    \eA{tell me, for I forget just now; but please come in and sit." "You're a lady
        of rank,"}
}
{
}

\Verse{88}
{
    \zA{请家里坐。”}
}
{
    \eA{answered goody Liu smiling,}
}
{
}

\Verse{89}
{
    \zA{刘老老一面走，一面笑说道：“你老是‘贵人多忘事’了， 那里还记得我们？”}
}
{
    \eA{as she walked along, "and do forget many things. How could you remember such
        as ourselves?" With these words still in her mouth, they had entered the
        house, whereupon Mrs. Chou ordered a hired waiting-maid to pour the tea.}
}
{
}

\Verse{90}
{
    \zA{说着，来至房中，周瑞家的命雇的小丫头倒上茶来吃着。}
}
{
    \eA{While they were having their tea she remarked, "How Pan Erh has managed to
        grow!" and then went on to make inquiries on the subject of various matters,}
}
{
}

\Verse{91}
{
    \zA{周瑞 家的又问道：“板儿长了这么大了么！”}
}
{
    \eA{which had occurred after their separation. "To-day," she also asked of goody
        Liu, "were you simply passing by?}
}
{
}

\Verse{92}
{
    \zA{又问些别后闲话。}
}
{
    \eA{or did you come with any express object?" "I've come, the fact is,}
}
{
}

\Verse{93}
{
    \zA{又问刘老老：“今日还 是路过，还是特来的？”}
}
{
    \eA{with an object!" promptly replied goody Liu; "(first of all) to see you, my
        dear sister-in-law; and, in the second place also,}
}
{
}

\Verse{94}
{
    \zA{刘老老便说：“原是特来瞧瞧嫂子；二则也请请姑太太的 安。}
}
{
    \eA{to inquire after my lady's health. If you could introduce me to see her for
        a while, it would be better; but if you can't, I must readily borrow your
        good offices,}
}
{
}

\Verse{95}
{
    \zA{若可以领我见一见更好，若不能，就借重嫂子转致意罢了。”}
}
{
    \eA{my sister-in-law, to convey my message." Mr. Chou Jui's wife, after
        listening to these words, at once became to a great extent aware of the
        object of her visit.}
}
{
}

\Verse{96}
{
    \zA{周瑞家的听了，便已猜着几分来意。}
}
{
    \eA{Her husband had, however, in years gone by in his attempt to purchase some
        land,}
}
{
}

\Verse{97}
{
    \zA{只因他丈夫昔年争买田地一事，多得狗儿 他父亲之力，今见刘老老如此，心中难却其意；二则也要显弄自己的体面。}
}
{
    \eA{obtained considerably the support of Kou Erh, so that when she, on this
        occasion, saw goody Liu in such a dilemma, she could not make up her mind to
        refuse her wish. Being in the second place keen upon making a display of her
        own respectability,}
}
{
}

\Verse{98}
{
    \zA{便笑说： “老老你放心。}
}
{
    \eA{she therefore said smilingly: "Old goody Liu, pray compose your mind!}
}
{
}

\Verse{99}
{
    \zA{大远的诚心诚意来了，岂有个不叫你见个真佛儿去的呢。}
}
{
    \eA{You've come from far off with a pure heart and honest purpose, and how can I
        ever not show you the way how to see this living Buddha?}
}
{
}

\Verse{100}
{
    \zA{论理，人 来客至，却都不与我相干。}
}
{
    \eA{Properly speaking, when people come and guests arrive, and verbal messages
        have to be given,}
}
{
}

\Verse{101}
{
    \zA{我们这里都是各一样儿：我们男的只管春秋两季地租子， 闲了时带着小爷们出门就完了；我只管跟太太奶奶们出门的事。}
}
{
    \eA{these matters are not any of my business, as we all here have each one kind
        of duties to carry out. My husband has the special charge of the rents of
        land coming in, during the two seasons of spring and autumn, and when at
        leisure, he takes the young gentlemen out of doors,}
}
{
}

\Verse{102}
{
    \zA{皆因你是太太的亲 戚，又拿我当个人，投奔了我来，我竟破个例给你通个信儿去。}
}
{
    \eA{and then his business is done. As for myself, I have to accompany my lady
        and young married ladies on anything connected with out-of-doors; but as you
        are a relative of my lady and have besides treated me as a high person and
        come to me for help,}
}
{
}

\Verse{103}
{
    \zA{但只一件，你还不 知道呢：我们这里不比五年前了。}
}
{
    \eA{I'll, after all, break this custom and deliver your message. There's only
        one thing, however, and which you, old lady,}
}
{
}

\Verse{104}
{
    \zA{如今太太不理事，都是琏二奶奶当家。}
}
{
    \eA{don't know. We here are not what we were five years before. My lady now
        doesn't much worry herself about anything;}
}
{
}

\Verse{105}
{
    \zA{你打量琏 二奶奶是谁？}
}
{
    \eA{and it's entirely lady Secunda who looks after the menage.}
}
{
}

\Verse{106}
{
    \zA{就是太太的内侄女儿，大舅老爷的女孩儿，小名儿叫凤哥的。”}
}
{
    \eA{But who do you presume is this lady Secunda? She's the niece of my lady, and
        the daughter of my master, the eldest maternal uncle of by-gone days.}
}
{
}

\Verse{107}
{
    \zA{刘老 老听了，忙问道：“原来是他？}
}
{
    \eA{Her infant name was Feng Ko." "Is it really she?" inquired promptly goody
        Liu,}
}
{
}

\Verse{108}
{
    \zA{怪道呢，我当日就说他不错。}
}
{
    \eA{after this explanation. "Isn't it strange? what I said about her years back
        has come out quite correct;}
}
{
}

\Verse{109}
{
    \zA{这么说起来，我今儿 还得见他了？”}
}
{
    \eA{but from all you say, shall I to-day be able to see her?" "That goes without
        saying,"}
}
{
}

\Verse{110}
{
    \zA{周瑞家的道：“这个自然。}
}
{
    \eA{replied Chou Jui's wife; "when any visitors come now-a-days,}
}
{
}

\Verse{111}
{
    \zA{如今有客来，都是凤姑娘周旋接待。}
}
{
    \eA{it's always lady Feng who does the honours and entertains them, and it's
        better to-day that you should see her for a while,}
}
{
}

\Verse{112}
{
    \zA{今 儿宁可不见太太，倒得见他一面，才不枉走这一遭儿。”}
}
{
    \eA{for then you will not have walked all this way to no purpose." "O mi to fu!"
        exclaimed old goody Liu; "I leave it entirely to your convenience,}
}
{
}

\Verse{113}
{
    \zA{刘老老道：“阿弥陀佛! 这全仗嫂子方便了。”}
}
{
    \eA{sister-in-law." "What's that you're saying?" observed Chou Jui's wife. "The
        proverb says: 'Our convenience is the convenience of others.'}
}
{
}

\Verse{114}
{
    \zA{周瑞家的说：“老老说那里话。}
}
{
    \eA{All I have to do is to just utter one word, and what trouble will that be to
        me."}
}
{
}

\Verse{115}
{
    \zA{俗语说的好：‘与人方便， 自己方便。’}
}
{
    \eA{Saying this, she bade the young waiting maid go to the side pavilion, and
        quietly ascertain whether,}
}
{
}

\Verse{116}
{
    \zA{不过用我一句话，又费不着我什么事。”}
}
{
    \eA{in her old ladyship's apartment, table had been laid. The young waiting-maid
        went on this errand,}
}
{
}

\Verse{117}
{
    \zA{说着，便唤小丫头：“到倒 厅儿上，悄悄的打听老太太屋里摆了饭了没有。”}
}
{
    \eA{and during this while, the two of them continued a conversation on certain
        irrelevant matters. "This lady Feng," observed goody Liu, "can this year be
        no older than twenty, and yet so talented as to manage such a household as
        this!}
}
{
}

\Verse{118}
{
    \zA{小丫头去了。}
}
{
    \eA{the like of her is not easy to find!"}
}
{
}

\Verse{119}
{
    \zA{这里二人又说了些闲话。}
}
{
    \eA{"Hai! my dear old goody," said Chou Jui's wife, after listening to her,}
}
{
}

\Verse{120}
{
    \zA{刘老老因说：“这位凤姑娘，今年不过十八九岁罢了， 就这等有本事，当这样的家，可是难得的！”}
}
{
    \eA{"it's not easy to explain; but this lady Feng, though young in years, is
        nevertheless, in the management of affairs, superior to any man. She has now
        excelled the others and developed the very features of a beautiful young
        woman.}
}
{
}

\Verse{121}
{
    \zA{周瑞家的听了道：“？}
}
{
    \eA{To say the least, she has ten thousand eyes in her heart,}
}
{
}

\Verse{122}
{
    \zA{!我的老老， 告诉不得你了!这凤姑娘年纪儿虽小，行事儿比是人都大呢。}
}
{
    \eA{and were they willing to wager their mouths, why ten men gifted with
        eloquence couldn't even outdo her! But by and bye, when you've seen her,
        you'll know all about her! There's only this thing,}
}
{
}

\Verse{123}
{
    \zA{如今出挑的美人儿似 的，少说着只怕有一万心眼子；再要赌口齿，十个会说的男人也说不过他呢。}
}
{
    \eA{she can't help being rather too severe in her treatment of those below her."
        While yet she spake, the young waiting-maid returned. "In her venerable
        lady's apartment," she reported, "repast has been spread, and already
        finished; lady Secunda is in madame Wang's chamber."}
}
{
}

\Verse{124}
{
    \zA{回来 你见了就知道了。}
}
{
    \eA{As soon as Chou Jui's wife heard this news, she speedily got up and pressed
        goody Liu to be off at once.}
}
{
}

\Verse{125}
{
    \zA{就只一件，待下人未免太严些儿。”}
}
{
    \eA{"This is," she urged, "just the hour for her meal, and as she is free we had
        better first go and wait for her;}
}
{
}

\Verse{126}
{
    \zA{说着，小丫头回来说：“老 太太屋里摆完了饭了，二奶奶在太太屋里呢。”}
}
{
    \eA{for were we to be even one step too late, a crowd of servants will come with
        their reports, and it will then be difficult to speak to her; and after her
        siesta, she'll have still less time to herself."}
}
{
}

\Verse{127}
{
    \zA{周瑞家的听了连忙起身，催着刘老 老：“快走，这一下来就只吃饭是个空儿，咱们先等着去。}
}
{
    \eA{As she passed these remarks, they all descended the couch together. Goody
        Liu adjusted their dresses, and, having impressed a few more words of advice
        on Pan Erh, they followed Chou Jui's wife through winding passages to Chia
        Lien's house.}
}
{
}

\Verse{128}
{
    \zA{若迟了一步，回事的人 多了，就难说了。}
}
{
    \eA{They came in the first instance into the side pavilion, where Chou Jui's
        wife placed old goody Liu to wait a little,}
}
{
}

\Verse{129}
{
    \zA{再歇了中觉，越发没时候了。”}
}
{
    \eA{while she herself went ahead, past the screen-wall and into the entrance of
        the court.}
}
{
}

\Verse{130}
{
    \zA{说着，一齐下了炕，整顿衣服， 又教了板儿几句话，跟着周瑞家的，逶迤往贾琏的住宅来。}
}
{
    \eA{Hearing that lady Feng had not come out, she went in search of an elderly
        waiting-maid of lady Feng, P'ing Erh by name, who enjoyed her confidence, to
        whom Chou Jui's wife first recounted from beginning to end the history of
        old goody Liu.}
}
{
}

\Verse{131}
{
    \zA{先至倒厅，周瑞家的将刘老老安插住等着，自己却先过影壁，走进了院门，知 凤姐尚未出来，先找着凤姐的一个心腹通房大丫头名唤平儿的。}
}
{
    \eA{"She has come to-day," she went on to explain, "from a distance to pay her
        obeisance. In days gone by, our lady used often to meet her, so that, on
        this occasion, she can't but receive her; and this is why I've brought her
        in!}
}
{
}

\Verse{132}
{
    \zA{周瑞家的先将刘老 老起初来历说明，又说：“今日大远的来请安，当日太太是常会的，所以我带了他 过来。}
}
{
    \eA{I'll wait here for lady Feng to come down, and explain everything to her;
        and I trust she'll not call me to task for officious rudeness." P'ing Erh,
        after hearing what she had to say, speedily devised the plan of asking them
        to walk in, and to sit there pending (lady Feng's arrival),}
}
{
}

\Verse{133}
{
    \zA{等着奶奶下来，我细细儿的回明了，想来奶奶也不至嗔着我莽撞的。”}
}
{
    \eA{when all would be right. Chou Jui's wife thereupon went out and led them in.
        When they ascended the steps of the main apartment, a young waiting-maid
        raised a red woollen portière,}
}
{
}

\Verse{134}
{
    \zA{平儿 听了，便作了个主意：“叫他们进来，先在这里坐着就是了。”}
}
{
    \eA{and as soon as they entered the hall, they smelt a whiff of perfume as it
        came wafted into their faces: what the scent was they could not
        discriminate; but their persons felt as if they were among the clouds.}
}
{
}

\Verse{135}
{
    \zA{周瑞家的才出去领 了他们进来，上了正房台阶，小丫头打起猩红毡帘，才入堂屋，只闻一阵香扑了脸 来，竟不知是何气味，身子就像在云端里一般。}
}
{
    \eA{The articles of furniture and ornaments in the whole room were all so
        brilliant to the sight, and so vying in splendour that they made the head to
        swim and the eyes to blink, and old goody Liu did nothing else the while
        than nod her head, smack her lips and invoke Buddha. Forthwith she was led
        to the eastern side into the suite of apartments, where was the bedroom of
        Chia Lien's eldest daughter.}
}
{
}

\Verse{136}
{
    \zA{满屋里的东西都是耀眼争光，使人 头晕目眩，刘老老此时只有点头咂嘴念佛而已。}
}
{
    \eA{P'ing Erh, who was standing by the edge of the stove-couch, cast a couple of
        glances at old goody Liu, and felt constrained to inquire how she was, and
        to press her to have a seat. Goody Liu, noticing that P'ing Erh was entirely
        robed in silks,}
}
{
}

\Verse{137}
{
    \zA{于是走到东边这间屋里，乃是贾琏 的女儿睡觉之所。}
}
{
    \eA{that she had gold pins fixed in her hair, and silver ornaments in her
        coiffure, and that her countenance resembled a flower or the moon (in
        beauty),}
}
{
}

\Verse{138}
{
    \zA{平儿站在炕沿边，打量了刘老老两眼，只得问个好，让了坐。}
}
{
    \eA{readily imagined her to be lady Feng, and was about to address her as my
        lady; but when she heard Mrs. Chou speak to her as Miss P'ing, and P'ing Erh
        promptly address Chou Jui's wife as Mrs. Chou,}
}
{
}

\Verse{139}
{
    \zA{刘 老老见平儿遍身绫罗，插金戴银，花容月貌，便当是凤姐儿了，才要称“姑奶奶”， 只见周瑞家的说：“他是平姑娘。”}
}
{
    \eA{she eventually became aware that she could be no more than a waiting-maid of
        a certain respectability. She at once pressed old goody Liu and Pan Erh to
        take a seat on the stove-couch. P'ing Erh and Chou Jui's wife sat face to
        face, on the edges of the couch. The waiting-maids brought the tea.}
}
{
}

\Verse{140}
{
    \zA{又见平儿赶着周瑞家的叫他“周大娘”，方知 不过是个有体面的丫头。}
}
{
    \eA{After they had partaken of it, old goody Liu could hear nothing but a "lo
        tang, lo tang" noise, resembling very much the sound of a bolting frame
        winnowing flour, and she could not resist looking now to the East,}
}
{
}

\Verse{141}
{
    \zA{于是让刘老老和板儿上了炕，平儿和周瑞家的对面坐在炕 沿上，小丫头们倒了茶来吃了。}
}
{
    \eA{and now to the West. Suddenly in the great Hall, she espied, suspended on a
        pillar, a box at the bottom of which hung something like the weight of a
        balance, which incessantly wagged to and fro. "What can this thing be?"
        communed goody Liu in her heart,}
}
{
}

\Verse{142}
{
    \zA{刘老老只听见咯当咯当的响声，很似打罗筛面的一般，不免东瞧西望的，忽见 堂屋中柱子上挂着一个匣子，底下又坠着一个秤铊似的，却不住的乱晃。}
}
{
    \eA{"What can be its use?" While she was aghast, she unexpectedly heard a sound
        of "tang" like the sound of a golden bell or copper cymbal, which gave her
        quite a start. In a twinkle of the eyes followed eight or nine consecutive
        strokes; and she was bent upon inquiring what it was, when she caught sight
        of several waiting-maids enter in a confused crowd.}
}
{
}

\Verse{143}
{
    \zA{刘老老心 中想着：“这是什么东西？}
}
{
    \eA{"Our lady has come down!" they announced. P'ìng Erh, together with Chou
        Jui's wife,}
}
{
}

\Verse{144}
{
    \zA{有煞用处呢？”}
}
{
    \eA{rose with all haste. "Old goody Liu,"}
}
{
}

\Verse{145}
{
    \zA{正发呆时，陡听得当的一声又若金钟铜 磬一般，倒吓得不住的展眼儿。}
}
{
    \eA{they urged, "do sit down and wait till it's time, when we'll come and ask
        you in." Saying this, they went out to meet lady Feng. Old goody Liu, with
        suppressed voice and ear intent, waited in perfect silence.}
}
{
}

\Verse{146}
{
    \zA{接着一连又是八九下，欲待问时，只见小丫头们一 齐乱跑，说：“奶奶下来了。”}
}
{
    \eA{She heard at a distance the voices of some people laughing, whereupon about
        ten or twenty women, with rustling clothes and petticoats, made their
        entrance, one by one, into the hall, and thence into the room on the other
        quarter.}
}
{
}

\Verse{147}
{
    \zA{平儿和周瑞家的忙起身说：“老老只管坐着，等是 时候儿我们来请你。”}
}
{
    \eA{She also detected two or three women, with red-lacquered boxes in their
        hands, come over on this part and remain in waiting. "Get the repast ready!"
        she heard some one from the offside say.}
}
{
}

\Verse{148}
{
    \zA{说着迎出去了。}
}
{
    \eA{The servants gradually dispersed and went out;}
}
{
}

\Verse{149}
{
    \zA{刘老老只屏声侧耳默候。}
}
{
    \eA{and there only remained in attendance a few of them to bring in the courses.}
}
{
}

\Verse{150}
{
    \zA{只听远远有人笑声， 约有一二十个妇人，衣裙？？}
}
{
    \eA{For a long time, not so much as the caw of a crow could be heard, when she
        unexpectedly perceived two servants carry in a couch-table,}
}
{
}

\Verse{151}
{
    \zA{，渐入堂屋，往那边屋内去了。}
}
{
    \eA{and lay it on this side of the divan. Upon this table were placed bowls and
        plates,}
}
{
}

\Verse{152}
{
    \zA{又见三两个妇人，都 捧着大红油漆盒进这边来等候。}
}
{
    \eA{in proper order replete, as usual, with fish and meats; but of these only a
        few kinds were slightly touched. As soon as Pan Erh perceived (all these
        delicacies),}
}
{
}

\Verse{153}
{
    \zA{听得那边说道“摆饭”，渐渐的人才散出去，只有 伺候端菜的几个人。}
}
{
    \eA{he set up such a noise, and would have some meat to eat, but goody Liu
        administered to him such a slap, that he had to keep away. Suddenly, she saw
        Mrs. Chou approach,}
}
{
}

\Verse{154}
{
    \zA{半日鸦雀不闻。}
}
{
    \eA{full of smiles, and as she waved her hand,}
}
{
}

\Verse{155}
{
    \zA{忽见两个人抬了一张炕桌来，放在这边炕上， 桌上碗盘摆列，仍是满满的鱼肉，不过略动了几样。}
}
{
    \eA{she called her. Goody Liu understood her meaning, and at once pulling Pan
        Erh off the couch, she proceeded to the centre of the Hall; and after Mrs.
        Chou had whispered to her again for a while, they came at length with slow
        step into the room on this side,}
}
{
}

\Verse{156}
{
    \zA{板儿一见就吵着要肉吃，刘老 老打了他一巴掌。}
}
{
    \eA{where they saw on the outside of the door, suspended by brass hooks, a deep
        red flowered soft portière. Below the window, on the southern side,}
}
{
}

\Verse{157}
{
    \zA{忽见周瑞家的笑嘻嘻走过来，点手儿叫他。}
}
{
    \eA{was a stove-couch, and on this couch was spread a crimson carpet. Leaning
        against the wooden partition wall, on the east side, stood a
        chain-embroidered back-cushion and a reclining pillow. There was also spread
        a large watered satin sitting cushion with a gold embroidered centre, and on
        the side stood cuspidores made of silver.}
}
{
}

\Verse{158}
{
    \zA{刘老老会意，于是带着板儿下炕。}
}
{
    \eA{Lady Feng, when at home, usually wore on her head a front-piece of dark
        martin à la Chao Chün, surrounded with tassels of strung pearls.}
}
{
}

\Verse{159}
{
    \zA{至堂屋中间，周瑞家的又和他咕唧了一会子，方蹭到这边屋内，只见门外铜钩上悬 着大红洒花软帘，南窗下是炕，炕上大红条毡，靠东边板壁立着一个锁子锦的靠背
        和一个引枕，铺着金线闪的大坐褥，傍边有银唾盒，那凤姐家常带着紫貂昭君套， 围着那攒珠勒子，穿着桃红洒花袄，石青刻丝灰鼠披风，大红洋绉银鼠皮裙，粉光
        脂艳，端端正正坐在那里，手内拿着小铜火箸儿拨手炉内的灰。}
}
{
    \eA{She had on a robe of peach-red flowered satin, a short pelisse of slate-blue
        stiff silk, lined with squirrel, and a jupe of deep red foreign crepe, lined
        with ermine. Resplendent with pearl-powder and with cosmetics, she sat in
        there, stately and majestic, with a small brass poker in her hands, with
        which she was stirring the ashes of the hand-stove.}
}
{
}

\Verse{160}
{
    \zA{平儿站在炕沿边， 捧着小小的一个填漆茶盘，盘内一个小盖钟儿。}
}
{
    \eA{P'ing Erh stood by the side of the couch, holding a very small lacquered
        tea-tray. In this tray was a small tea-cup with a cover. Lady Feng neither
        took any tea, nor did she raise her head,}
}
{
}

\Verse{161}
{
    \zA{凤姐也不接茶，也不抬头，只管拨 那灰，慢慢的道：“怎么还不请进来？”}
}
{
    \eA{but was intent upon stirring the ashes of the hand-stove. "How is it you
        haven't yet asked her to come in?" she slowly inquired; and as she spake,
        she turned herself round and was about to ask for some tea,}
}
{
}

\Verse{162}
{
    \zA{一面说，一面抬身要茶时，只见周瑞家的 已带了两个人立在面前了，这才忙欲起身、犹未起身，满面春风的问好，又嗔着周 瑞家的：“怎么不早说！”}
}
{
    \eA{when she perceived that Mrs. Chou had already introduced the two persons and
        that they were standing in front of her. She forthwith pretended to rise,
        but did not actually get up, and with a face radiant with smiles, she
        ascertained about their health, after which she went in to chide Chou Jui's
        wife. "Why didn't you tell me they had come before?" she said. Old goody Liu
        was already by this time prostrated on the ground,}
}
{
}

\Verse{163}
{
    \zA{刘老老已在地下拜了几拜，问姑奶奶安。}
}
{
    \eA{and after making several obeisances, "How are you, my lady?" she inquired.
        "Dear Mrs. Chou," lady Feng immediately observed,}
}
{
}

\Verse{164}
{
    \zA{凤姐忙说：“周 姐姐，搀着不拜罢。}
}
{
    \eA{"do pull her up, and don't let her prostrate herself! I'm yet young in years
        and don't know her much; what's more, I've no idea what's the degree of the
        relationship between us, and I daren't speak directly to her."}
}
{
}

\Verse{165}
{
    \zA{我年轻，不大认得，可也不知是什么辈数儿，不敢称呼。”}
}
{
    \eA{"This is the old lady about whom I spoke a short while back," speedily
        explained Mrs. Chou. Lady Feng nodded her head assentingly. By this time old
        goody Liu had taken a seat on the edge of the stove-couch.}
}
{
}

\Verse{166}
{
    \zA{周 瑞家的忙回道：“这就是我才回的那个老老了。”}
}
{
    \eA{As for Pan Erh, he had gone further, and taken refuge behind her back; and
        though she tried, by every means, to coax him to come forward and make a
        bow,}
}
{
}

\Verse{167}
{
    \zA{凤姐点头，刘老老已在炕沿上坐 下了，板儿便躲在他背后，百般的哄他出来作揖，他死也不肯。}
}
{
    \eA{he would not, for the life of him, consent. "Relatives though we be,"
        remarked lady Feng, as she smiled, "we haven't seen much of each other, so
        that our relations have been quite distant. But those who know how matters
        stand will assert that you all despise us,}
}
{
}

\Verse{168}
{
    \zA{凤姐笑道：“亲戚们不大走动，都疏远了。}
}
{
    \eA{and won't often come to look us up; while those mean people, who don't know
        the truth, will imagine that we have no eyes to look at any one."}
}
{
}

\Verse{169}
{
    \zA{知道的呢，说你们弃嫌我们，不肯 常来，不知道的那起小人，还只当我们眼里没人似的。”}
}
{
    \eA{Old goody Liu promptly invoked Buddha. "We are at home in great straits,"
        she pleaded, "and that's why it wasn't easy for us to manage to get away and
        come! Even supposing we had come as far as this, had we not given your
        ladyship a slap on the mouth,}
}
{
}

\Verse{170}
{
    \zA{刘老老忙念佛道：“我们 家道艰难，走不起。}
}
{
    \eA{those gentlemen would also, in point of fact, have looked down upon us as a
        mean lot." "Why, language such as this,"}
}
{
}

\Verse{171}
{
    \zA{来到这里，没的给姑奶奶打嘴，就是管家爷们瞧着也不像。”}
}
{
    \eA{exclaimed lady Feng smilingly, "cannot help making one's heart full of
        displeasure! We simply rely upon the reputation of our grandfather to
        maintain the status of a penniless official; that's all! Why, in whose
        household is there anything substantial?}
}
{
}

\Verse{172}
{
    \zA{凤姐笑道：“这话没的叫人恶心。}
}
{
    \eA{we are merely the denuded skeleton of what we were in days of old, and no
        more! As the proverb has it:}
}
{
}

\Verse{173}
{
    \zA{不过托赖着祖父的虚名，作个穷官儿罢咧，谁家 有什么？}
}
{
    \eA{The Emperor himself has three families of poverty-stricken relatives; and
        how much more such as you and I?" Having passed these remarks,}
}
{
}

\Verse{174}
{
    \zA{不过也是个空架子，俗语儿说的好，‘朝廷还有三门子穷亲’呢，何况你 我。”}
}
{
    \eA{she inquired of Mrs. Chou, "Have you let madame know, yes or no?" "We are
        now waiting," replied Mrs. Chou, "for my lady's orders." "Go and have a
        look," said lady Feng;}
}
{
}

\Verse{175}
{
    \zA{说着，又问周瑞家的：“回了太太了没有？”}
}
{
    \eA{"but, should there be any one there, or should she be busy, then don't make
        any mention; but wait until she's free,}
}
{
}

\Verse{176}
{
    \zA{周瑞家的道：“等奶奶的示下。”}
}
{
    \eA{when you can tell her about it and see what she says." Chou Jui's wife,
        having expressed her compliance,}
}
{
}

\Verse{177}
{
    \zA{凤姐儿道：“你去瞧瞧，要是有人就罢；要得闲呢，就回了，看怎么说。”}
}
{
    \eA{went off on this errand. During her absence, lady Feng gave orders to some
        servants to take a few fruits and hand them to Pan Erh to eat; and she was
        inquiring about one thing and another,}
}
{
}

\Verse{178}
{
    \zA{周瑞家 的答应去了。}
}
{
    \eA{when there came a large number of married women,}
}
{
}

\Verse{179}
{
    \zA{这里凤姐叫人抓了些果子给板儿吃，刚问了几句闲话时，就有家下许多媳妇儿 管事的来回话。}
}
{
    \eA{who had the direction of affairs in the household, to make their several
        reports. P'ing Erh announced their arrival to lady Feng, who said: "I'm now
        engaged in entertaining some guests, so let them come back again in the
        evening;}
}
{
}

\Verse{180}
{
    \zA{平儿回了，凤姐道：“我这里陪客呢，晚上再来回。}
}
{
    \eA{but should there be anything pressing then bring it in and I'll settle it at
        once." P'ing Erh left the room, but she returned in a short while.}
}
{
}

\Verse{181}
{
    \zA{要有紧事，你 就带进来现办。”}
}
{
    \eA{"I've asked them," she observed, "but as there's nothing of any urgency,}
}
{
}

\Verse{182}
{
    \zA{平儿出去，一会进来说：“我问了，没什么要紧的。}
}
{
    \eA{I told them to disperse." Lady Feng nodded her head in token of approval,
        when she perceived Chou Jui's wife come back. "Our lady," she reported,}
}
{
}

\Verse{183}
{
    \zA{我叫他们散 了。”}
}
{
    \eA{as she addressed lady Feng, "says that she has no leisure to-day,}
}
{
}

\Verse{184}
{
    \zA{凤姐点头。}
}
{
    \eA{that if you, lady Secunda,}
}
{
}

\Verse{185}
{
    \zA{只见周瑞家的回来，向凤姐道：“太太说：‘今日不得闲儿，二 奶奶陪着也是一样，多谢费心想着。}
}
{
    \eA{will entertain them, it will come to the same thing; that she's much obliged
        for their kind attention in going to the trouble of coming; that if they
        have come simply on a stroll, then well and good, but that if they have
        aught to say, they should tell you, lady Secunda,}
}
{
}

\Verse{186}
{
    \zA{要是白来逛逛呢便罢；有什么说的，只管告诉 二奶奶。’”}
}
{
    \eA{which will be tantamount to their telling her." "I've nothing to say,"
        interposed old goody Liu. "I simply come to see our elder and our younger
        lady, which is a duty on my part, a relative as I am." "Well,}
}
{
}

\Verse{187}
{
    \zA{刘老老道：“也没甚的说，不过来瞧瞧姑太太姑奶奶，也是亲戚们的 情分。”}
}
{
    \eA{if there's nothing particular that you've got to say, all right," Mrs. Chou
        forthwith added, "but if you do have anything, don't hesitate telling lady
        Secunda, and it will be just as if you had told our lady."}
}
{
}

\Verse{188}
{
    \zA{周瑞家的道：“没有什么说的便罢；要有话，只管回二奶奶，和太太是一 样儿的。”}
}
{
    \eA{As she uttered these words, she winked at goody Liu. Goody Liu understood
        what she meant, but before she could give vent to a word, her face got
        scarlet, and though she would have liked not to make any mention of the
        object of her visit,}
}
{
}

\Verse{189}
{
    \zA{一面说一面递了个眼色儿。}
}
{
    \eA{she felt constrained to suppress her shame and to speak out. "Properly
        speaking," she observed,}
}
{
}

\Verse{190}
{
    \zA{刘老老会意，未语先红了脸。}
}
{
    \eA{"this being the first time I see you, my lady, I shouldn't mention what I've
        to say,}
}
{
}

\Verse{191}
{
    \zA{待要不说，今 日所为何来？}
}
{
    \eA{but as I come here from far off to seek your assistance, my old friend, I
        have no help but to mention it."}
}
{
}

\Verse{192}
{
    \zA{只得勉强说道：“论今日初次见，原不该说的，只是大远的奔了你老 这里来，少不得说了……”刚说到这里，只听二门上小厮们回说：“东府里小大爷 进来了。”}
}
{
    \eA{She had barely spoken as much as this, when she heard the youths at the
        inner-door cry out: "The young gentleman from the Eastern Mansion has come."
        Lady Feng promptly interrupted her. "Old goody Liu," she remarked, "you
        needn't add anything more."}
}
{
}

\Verse{193}
{
    \zA{凤姐忙和刘老老摆手道：“不必说了。”}
}
{
    \eA{She, at the same time, inquired, "Where's your master, Mr. Jung?" when
        became audible the sound of footsteps along the way,}
}
{
}

\Verse{194}
{
    \zA{一面便问：“你蓉大爷在那里 呢？”}
}
{
    \eA{and in walked a young man of seventeen or eighteen. His appearance was
        handsome, his person slender and graceful.}
}
{
}

\Verse{195}
{
    \zA{只听一路靴子响，进来了一个十七八岁的少年，面目清秀，身段苗条，美服 华冠，轻裘宝带。}
}
{
    \eA{He had on light furs, a girdle of value, costly clothes and a beautiful cap.
        At this stage, goody Liu did not know whether it was best to sit down or to
        stand up, neither could she find anywhere to hide herself. "Pray sit down,"
        urged lady Feng,}
}
{
}

\Verse{196}
{
    \zA{刘老老此时坐不是站不是，藏没处藏，躲没处躲。}
}
{
    \eA{with a laugh; "this is my nephew!' Old goody Liu then wriggled herself, now
        one way, and then another, on to the edge of the couch,}
}
{
}

\Verse{197}
{
    \zA{凤姐笑道：“你 只管坐着罢，这是我侄儿。”}
}
{
    \eA{where she took a seat. "My father," Chia Jung smilingly ventured, "has sent
        me to ask a favour of you, aunt. On some previous occasion,}
}
{
}

\Verse{198}
{
    \zA{刘老老才扭扭捏捏的在炕沿儿上侧身坐下。}
}
{
    \eA{our grand aunt gave you, dear aunt, a stove-couch glass screen, and as
        to-morrow father has invited some guests of high standing,}
}
{
}

\Verse{199}
{
    \zA{那贾蓉请了安，笑回道：“我父亲打发来求婶子，上回老舅太太给婶子的那架 玻璃炕屏，明儿请个要紧的客，略摆一摆就送来。”}
}
{
    \eA{he wishes to borrow it to lay it out for a little show; after which he
        purposes sending it back again." "You're late by a day," replied lady Feng.
        "It was only yesterday that I gave it to some one." Chia Jung, upon hearing
        this, forthwith, with giggles and smiles, made, near the edge of the couch,}
}
{
}

\Verse{200}
{
    \zA{凤姐道：“你来迟了，昨儿已 经给了人了。”}
}
{
    \eA{a sort of genuflexion. "Aunt," he went on, "if you don't lend it, father
        will again say that I don't know how to speak,}
}
{
}

\Verse{201}
{
    \zA{贾蓉听说，便笑嘻嘻的在炕沿上下个半跪道：“婶子要不借，我父 亲又说我不会说话了，又要挨一顿好打。}
}
{
    \eA{and I shall get another sound thrashing. You must have pity upon your
        nephew, aunt." "I've never seen anything like this," observed lady Feng
        sneeringly; "the things belonging to the Wang family are all good, but where
        have you put all those things of yours?}
}
{
}

\Verse{202}
{
    \zA{好婶子，只当可怜我罢！”}
}
{
    \eA{the only good way is that you shouldn't see anything of ours, for as soon as
        you catch sight of anything,}
}
{
}

\Verse{203}
{
    \zA{凤姐笑道：“也 没见我们王家的东西都是好的？}
}
{
    \eA{you at once entertain a wish to carry it off." "Pray, aunt," entreated Chia
        Jung with a smile, "do show me some compassion."}
}
{
}

\Verse{204}
{
    \zA{你们那里放着那些好东西，只别看见我的东西才 罢，一见了就想拿了去。”}
}
{
    \eA{"Mind your skin!" lady Feng warned him, "if you do chip or spoil it in the
        least." She then bade P'ing Erh take the keys of the door of the upstairs
        room and send for several trustworthy persons to carry it away.}
}
{
}

\Verse{205}
{
    \zA{贾蓉笑道：“只求婶娘开恩罢！”}
}
{
    \eA{Chia Jung was so elated that his eyebrows dilated and his eyes smiled.}
}
{
}

\Verse{206}
{
    \zA{凤姐道：“碰坏一点 儿，你可仔细你的皮！”}
}
{
    \eA{"I've brought myself," he added, with vehemence, "some men to take it away;
        I won't let them recklessly bump it about."}
}
{
}

\Verse{207}
{
    \zA{因命平儿拿了楼门上钥匙，叫几个妥当人来抬去。}
}
{
    \eA{Saying this, he speedily got up and left the room. Lady Feng suddenly
        bethought herself of something, and turning towards the window,}
}
{
}

\Verse{208}
{
    \zA{贾蓉喜 的眉开眼笑，忙说：“我亲自带人拿去，别叫他们乱碰。”}
}
{
    \eA{she called out, "Jung Erh, come back." Several servants who stood outside
        caught up her words: "Mr. Jung," they cried, "you're requested to go back;"}
}
{
}

\Verse{209}
{
    \zA{说着便起身出去了。}
}
{
    \eA{whereupon Chia Jung turned round and retraced his steps;}
}
{
}

\Verse{210}
{
    \zA{这 凤姐忽然想起一件事来，便向窗外叫：“蓉儿回来！”}
}
{
    \eA{and with hands drooping respectfully against his sides, he stood ready to
        listen to his aunt's wishes. Lady Feng was however intent upon gently
        sipping her tea,}
}
{
}

\Verse{211}
{
    \zA{外面几个人接声说：“请蓉 大爷回来呢！”}
}
{
    \eA{and after a good long while of abstraction, she at last smiled: "Never
        mind," she remarked; "you can go.}
}
{
}

\Verse{212}
{
    \zA{贾蓉忙回来，满脸笑容的瞅着凤姐，听何指示。}
}
{
    \eA{But come after you've had your evening meal, and I'll then tell you about
        it. Just now there are visitors here;}
}
{
}

\Verse{213}
{
    \zA{那凤姐只管慢慢吃 茶，出了半日神，忽然把脸一红，笑道：“罢了，你先去罢。}
}
{
    \eA{and besides, I don't feel in the humour." Chia Jung thereupon retired with
        gentle step. Old goody Liu, by this time, felt more composed in body and
        heart. "I've to-day brought your nephew,"}
}
{
}

\Verse{214}
{
    \zA{晚饭后你来再说罢。}
}
{
    \eA{she then explained, "not for anything else,}
}
{
}

\Verse{215}
{
    \zA{这会子有人，我也没精神了。”}
}
{
    \eA{but because his father and mother haven't at home so much as anything to
        eat; the weather besides is already cold,}
}
{
}

\Verse{216}
{
    \zA{贾蓉答应个是，抿着嘴儿一笑，方慢慢退去。}
}
{
    \eA{so that I had no help but to take your nephew along and come to you, old
        friend, for assistance!" As she uttered these words,}
}
{
}

\Verse{217}
{
    \zA{这刘老老方安顿了，便说道：“我今日带了你侄儿，不为别的，因他爹娘连吃 的没有，天气又冷，只得带了你侄儿奔了你老来。”}
}
{
    \eA{she again pushed Pan Erh forward. "What did your father at home tell you to
        say?" she asked of him; "and what did he send us over here to do? Was it
        only to give our minds to eating fruit?" Lady Feng had long ago understood
        what she meant to convey, and finding that she had no idea how to express
        herself in a decent manner,}
}
{
}

\Verse{218}
{
    \zA{说着，又推板儿道：“你爹在 家里怎么教你的？}
}
{
    \eA{she readily interrupted her with a smile. "You needn't mention anything,"
        she observed, "I'm well aware of how things stand;"}
}
{
}

\Verse{219}
{
    \zA{打发咱们来作煞事的？}
}
{
    \eA{and addressing herself to Mrs. Chou, she inquired, "Has this old lady had
        breakfast,}
}
{
}

\Verse{220}
{
    \zA{只顾吃果子！”}
}
{
    \eA{yes or no?" Old goody Liu hurried to explain.}
}
{
}

\Verse{221}
{
    \zA{凤姐早已明白了，听他不会 说话，因笑道：“不必说了，我知道了。”}
}
{
    \eA{"As soon as it was daylight," she proceeded, "we started with all speed on
        our way here, and had we even so much as time to have any breakfast?"}
}
{
}

\Verse{222}
{
    \zA{因问周瑞家的道：“这老老不知用了早 饭没有呢？”}
}
{
    \eA{Lady Feng promptly gave orders to send for something to eat. In a short
        while Chou Jui's wife had called for a table of viands for the guests,}
}
{
}

\Verse{223}
{
    \zA{刘老老忙道：“一早就往这里赶咧，那里还有吃饭的工夫咧？”}
}
{
    \eA{which was laid in the room on the eastern side, and then came to take goody
        Liu and Pan Erh over to have their repast. "My dear Mrs. Chou," enjoined
        lady Feng,}
}
{
}

\Verse{224}
{
    \zA{凤姐 便命快传饭来。}
}
{
    \eA{"give them all they want, as I can't attend to them myself;"}
}
{
}

\Verse{225}
{
    \zA{一时周瑞家的传了一桌客馔，摆在东屋里，过来带了刘老老和板儿 过去吃饭。}
}
{
    \eA{which said, they hastily passed over into the room on the eastern side. Lady
        Feng having again called Mrs. Chou, asked her: "When you first informed
        madame about them, what did she say?" "Our Lady observed,"}
}
{
}

\Verse{226}
{
    \zA{凤姐这里道：“周姐姐好生让着些儿，我不能陪了。”}
}
{
    \eA{replied Chou Jui's wife, "that they don't really belong to the same family;
        that, in former years, their grandfather was an official at the same place
        as our old master;}
}
{
}

\Verse{227}
{
    \zA{一面又叫过周瑞 家的来问道：“方才回了太太，太太怎么说了？”}
}
{
    \eA{that hence it came that they joined ancestors; that these few years there
        hasn't been much intercourse (between their family and ours); that some
        years back, whenever they came on a visit,}
}
{
}

\Verse{228}
{
    \zA{周瑞家的道：“太太说：‘他们 原不是一家子；当年他们的祖和太老爷在一处做官，因连了宗的。}
}
{
    \eA{they were never permitted to go empty-handed, and that as their coming on
        this occasion to see us is also a kind attention on their part, they
        shouldn't be slighted. If they've anything to say,"}
}
{
}

\Verse{229}
{
    \zA{这几年不大走动。}
}
{
    \eA{(our lady continued), "tell lady Secunda to do the necessary,}
}
{
}

\Verse{230}
{
    \zA{当时他们来了，却也从没空过的。}
}
{
    \eA{and that will be right." "Isn't it strange!" exclaimed lady Feng, as soon as
        she had heard the message;}
}
{
}

\Verse{231}
{
    \zA{如今来瞧我们，也是他的好意，别简慢了他。}
}
{
    \eA{"since we are all one family, how is it I'm not familiar even with so much
        as their shadow?" While she was uttering these words,}
}
{
}

\Verse{232}
{
    \zA{要 有什么话，叫二奶奶裁夺着就是了。’”}
}
{
    \eA{old goody Liu had had her repast and come over, dragging Pan Erh; and,
        licking her lips and smacking her mouth,}
}
{
}

\Verse{233}
{
    \zA{凤姐听了说道：“怪道既是一家子，我怎 么连影儿也不知道！”}
}
{
    \eA{she expressed her thanks. Lady Feng smiled. "Do pray sit down," she said,
        "and listen to what I'm going to tell you. What you, old lady, meant a
        little while back to convey,}
}
{
}

\Verse{234}
{
    \zA{说话间，刘老老已吃完了饭，拉了板儿过来，舔唇咂嘴的道谢。}
}
{
    \eA{I'm already as much as yourself well acquainted with! Relatives, as we are,
        we shouldn't in fact have waited until you came to the threshold of our
        doors, but ought, as is but right,}
}
{
}

\Verse{235}
{
    \zA{凤姐笑道：“且 请坐下，听我告诉你：方才你的意思，我已经知道了。}
}
{
    \eA{to have attended to your needs. But the thing is that, of late, the
        household affairs are exceedingly numerous, and our lady, advanced in years
        as she is,}
}
{
}

\Verse{236}
{
    \zA{论起亲戚来，原该不等上门 就有照应才是；但只如今家里事情太多，太太上了年纪，一时想不到是有的。}
}
{
    \eA{couldn't at a moment, it may possibly be, bethink herself of you all! What's
        more, when I took over charge of the management of the menage, I myself
        didn't know of all these family connections! Besides, though to look at us
        from outside everything has a grand and splendid aspect,}
}
{
}

\Verse{237}
{
    \zA{我如 今接着管事，这些亲戚们又都不大知道，况且外面看着虽是烈烈轰轰，不知大有大 的难处，说给人也未必信。}
}
{
    \eA{people aren't aware that large establishments have such great hardships,
        which, were we to recount to others, they would hardly like to credit as
        true. But since you've now come from a great distance, and this is the first
        occasion that you open your mouth to address me,}
}
{
}

\Verse{238}
{
    \zA{你既大远的来了，又是头一遭儿和我张个口，怎么叫你 空回去呢？}
}
{
    \eA{how can I very well allow you to return to your home with empty hands! By a
        lucky coincidence our lady gave, yesterday, to the waiting-maids, twenty
        taels to make clothes with,}
}
{
}

\Verse{239}
{
    \zA{可巧昨儿太太给我的丫头们作衣裳的二十两银子还没动呢，你不嫌少， 先拿了去用罢。”}
}
{
    \eA{a sum which they haven't as yet touched, and if you don't despise it as too
        little, you may take it home as a first instalment, and employ it for your
        wants." When old goody Liu heard the mention made by lady Feng of their
        hardships,}
}
{
}

\Verse{240}
{
    \zA{那刘老老先听见告艰苦，只当是没想头了；又听见给他二十两银 子，喜的眉开眼笑道：“我们也知道艰难的，但只俗语说的：‘瘦死的骆驼比马还 大’呢。}
}
{
    \eA{she imagined that there was no hope; but upon hearing her again speak of
        giving her twenty taels, she was exceedingly delighted, so much so that her
        eyebrows dilated and her eyes gleamed with smiles. "We too know," she
        smilingly remarked,}
}
{
}

\Verse{241}
{
    \zA{凭他怎样，你老拔一根寒毛比我们的腰还壮哩。”}
}
{
    \eA{"all about difficulties! but the proverb says, 'A camel dying of leanness is
        even bigger by much than a horse!' No matter what those distresses may be,}
}
{
}

\Verse{242}
{
    \zA{周瑞家的在旁听见他说 的粗鄙，只管使眼色止他。}
}
{
    \eA{were you yet to pluck one single hair from your body, my old friend, it
        would be stouter than our own waist." Chou Jui's wife stood by,}
}
{
}

\Verse{243}
{
    \zA{凤姐笑而不睬，叫平儿把昨儿那包银子拿来，再拿一串 钱，都送至刘老老跟前。}
}
{
    \eA{and on hearing her make these coarse utterances, she did all she could to
        give her a hint by winking, and make her desist. Lady Feng laughed and paid
        no heed; but calling P'ing Erh, she bade her fetch the parcel of money,}
}
{
}

\Verse{244}
{
    \zA{凤姐道：“这是二十两银子，暂且给这孩子们作件冬衣罢。}
}
{
    \eA{which had been given to them the previous day, and to also bring a string of
        cash; and when these had been placed before goody Liu's eyes: "This is,"
        said lady Feng,}
}
{
}

\Verse{245}
{
    \zA{改日没事，只管来逛逛，才是亲戚们的意思。}
}
{
    \eA{"silver to the amount of twenty taels, which was for the time given to these
        young girls to make winter clothes with;}
}
{
}

\Verse{246}
{
    \zA{天也晚了，不虚留你们了，到家该问 好的都问个好儿罢。”}
}
{
    \eA{but some other day, when you've nothing to do, come again on a stroll, in
        evidence of the good feeling which should exist between relatives. It's
        besides already late,}
}
{
}

\Verse{247}
{
    \zA{一面说，一面就站起来了。}
}
{
    \eA{and I don't wish to detain you longer and all for no purpose; but, on your
        return home,}
}
{
}

\Verse{248}
{
    \zA{刘老老只是千恩万谢的，拿了银钱，跟着周瑞家的走到外边。}
}
{
    \eA{present my compliments to all those of yours to whom I should send them." As
        she spake, she stood up. Old goody Liu gave utterance to a thousand and ten
        thousand expressions of gratitude,}
}
{
}

\Verse{249}
{
    \zA{周瑞家的道：“我 的娘!你怎么见了他倒不会说话了呢？}
}
{
    \eA{and taking the silver and cash, she followed Chou Jui's wife on her way to
        the out-houses. "Well, mother dear," inquired Mrs. Chou,}
}
{
}

\Verse{250}
{
    \zA{开口就是‘你侄儿’。}
}
{
    \eA{"what did you think of my lady that you couldn't speak;}
}
{
}

\Verse{251}
{
    \zA{我说句不怕你恼的话： 就是亲侄儿也要说的和软些儿。}
}
{
    \eA{and that whenever you opened your mouth it was all 'your nephew.' I'll make
        just one remark, and I don't mind if you do get angry. Had he even been your
        kindred nephew,}
}
{
}

\Verse{252}
{
    \zA{那蓉大爷才是他的侄儿呢。}
}
{
    \eA{you should in fact have been somewhat milder in your language; for that
        gentleman, Mr. Jung,}
}
{
}

\Verse{253}
{
    \zA{他怎么又跑出这么个侄 儿来了呢！”}
}
{
    \eA{is her kith and kin nephew, and whence has appeared such another nephew of
        hers (as Pan Erh)?"}
}
{
}

\Verse{254}
{
    \zA{刘老老笑道：“我的嫂子!我见了他，心眼儿里爱还爱不过来，那里 还说的上话来？”}
}
{
    \eA{Old goody Liu smiled. "My dear sister-in-law," she replied, "as I gazed upon
        her, were my heart and eyes, pray, full of admiration or not? and how then
        could I speak as I should?"}
}
{
}

\Verse{255}
{
    \zA{二人说着，又到周瑞家坐了片刻。}
}
{
    \eA{As they were chatting, they reached Chou Jui's house. They had been sitting
        for a while,}
}
{
}

\Verse{256}
{
    \zA{刘老老要留下一块银子给周家 的孩子们买果子吃，周瑞家的那里放在眼里，执意不肯。}
}
{
    \eA{when old goody Liu produced a piece of silver, which she was purposing to
        leave behind, to be given to the young servants in Chou Jui's house to
        purchase fruit to eat; but how could Mrs. Chou satiate her eye with such a
        small piece of silver?}
}
{
}

\Verse{257}
{
    \zA{刘老老感谢不尽，仍从后 门去了。}
}
{
    \eA{She was determined in her refusal to accept it, so that old goody Liu, after
        assuring her of her boundless gratitude,}
}
{
}

\Verse{258}
{
    \zA{未知去后如何，且听下回分解。}
}
{
    \eA{took her departure out of the back gate she had come in from. Reader, you do
        not know what happened after old goody Liu left,}
}
{
}

\Verse{259}
{
    \zA{(本章完)}
}
{
    \eA{but listen to the explanation which will be given in the next chapter.}
}
{
}
